\chapter{Conclusion}\label{chap:conclusion}
This thesis investigated the construction of a simulation-oriented digital twin for railway network simulation using SUMO 
and heterogeneous railway datasets. The objective was not to build a complete real-time digital twin, but to study the 
foundations required to transform infrastructure and operational data into executable simulation models.

The work was structured around two complementary levels of representation. The macro-level model represents the railway 
network as a station-to-station graph and focuses on the integration of real operational data. The micro-level model 
represents the infrastructure in more detail through tracks, platforms, switches, and routing structures. Together, these 
two approaches make it possible to evaluate the trade-off between simplicity, scalability, infrastructure realism, and 
simulation feasibility.

\section{Summary of Contributions}
The first contribution of this thesis is the design and implementation of a data processing pipeline able to transform 
railway datasets into SUMO-compatible artefacts. This includes the extraction and preprocessing of infrastructure and 
operational data, the generation of simulation networks, the construction of train routes, and the conversion of simulation 
outputs into analysable datasets.

The second contribution is the development of a macro-level railway simulation model. This model abstracts the railway 
network as stations connected by railway links. Although simplified, it enables the integration of real punctuality data 
and validates the complete workflow from raw data to simulation execution.

The third contribution is the development of a micro-level railway infrastructure model. This model introduces more 
detailed infrastructure processing steps, including switch detection, platform matching, track segmentation, and dual graph 
construction. It provides a richer representation of the railway network and makes it possible to test train circulation 
on a more detailed infrastructure.

Finally, this work contributes by identifying the main technical challenges involved in building railway digital twins from 
real-world data. These challenges include data alignment, infrastructure completeness, platform assignment, route 
reconstruction, SUMO-specific constraints, and computational scalability.

\section{Key Findings}
Regarding \textbf{RQ1}, the results show that heterogeneous railway infrastructure and operational datasets can be 
transformed into SUMO-compatible simulation components through a structured pipeline. The generated datasets provide a 
strong basis for simulation, especially at the station and platform level.

Regarding \textbf{RQ2}, the macro-level model proved to be the most robust representation for integrating real operational 
data. Since train movements can be reconstructed as station-to-station trajectories, the macro model supports real-data 
simulation on selected chains of stations. It therefore provides an effective first validation of the digital twin 
workflow.

Regarding \textbf{RQ3}, the micro-level model shows that detailed infrastructure reconstruction is feasible and useful. The 
generated network can represent tracks, switches, platforms, and segmented railway sections. Random-trip simulations 
confirm that the infrastructure can support train circulation in SUMO.\@

Regarding \textbf{RQ4}, the main limitation appears when real operational data is integrated into the micro-level model. 
Real train services require continuous and valid paths between consecutive stations. In the current implementation, this 
condition is not always satisfied because the reconstructed dual graph does not always provide valid paths between station 
pairs. This result is important because it shows that high dataset coverage does not necessarily imply operational 
validity. A railway digital twin must therefore be evaluated not only by the number of reconstructed elements, but also by 
its ability to support realistic train routes.

Overall, the main finding is that macro-level and micro-level models are complementary. The macro-level model is more 
robust and scalable for real-data simulation, while the micro-level model provides a richer basis for detailed 
infrastructure analysis but requires stronger topological consistency.

\section{Practical Implications}
The proposed framework shows that existing railway data can be transformed into executable simulation scenarios without 
manually building the entire network. This is relevant for researchers and railway stakeholders interested in 
simulation-based analysis, data integration, or early-stage digital twin development.

The macro-level model can be used for preliminary traffic analysis, selected line simulation, and validation of operational 
data integration. Its simplified structure makes it easier to execute and less sensitive to missing detailed infrastructure 
information.

The micro-level model is more limited in its current form, but it provides a foundation for future analyses involving 
platforms, switches, local bottlenecks, routing conflicts, and infrastructure capacity. Its main value lies in showing 
where the transition from a structural infrastructure model to an operational simulation model becomes difficult.

The work also confirms that SUMO can be used as a flexible platform for railway digital twin prototyping. However, several 
railway-specific elements, such as signalling, interlocking, dispatching, platform allocation, and conflict management, 
require additional modelling beyond the standard SUMO workflow.

\section{Future directions}
The work presented in this thesis provides a first implementation of a railway digital twin pipeline based on real 
infrastructure and operational data. Future work should focus both on improving the existing implementation and on 
extending the framework toward more advanced decision-support functionalities.

\subsection{Improving Pathfinding and Micro-Level Routing}
The first priority is to improve the micro-level pathfinding process. The current implementation relies on a dual graph 
representation, where track segments are treated as graph nodes and transitions between tracks define the possible 
movements. This approach is useful, but the results show that it is not sufficient in its current form to reconstruct real 
train services reliably.

Future work should therefore investigate alternative or complementary routing methods. One possibility would be to 
construct a directed infrastructure graph closer to SUMO’s internal edge representation and perform pathfinding directly on 
this graph. Another possibility would be to combine the dual graph with additional validation rules based on track 
direction, switch connectivity, platform accessibility, and known station-to-station connections. A hybrid routing strategy 
could also be considered, where macro-level paths are first used to constrain the search space before computing detailed 
micro-level routes.

Improving this part of the code would also require better diagnostics. Instead of only detecting that no path exists, the 
pipeline should identify why the path failed: missing edge, wrong direction, disconnected platform, switch inconsistency, 
or incomplete station mapping. This would make debugging easier and would help transform routing failures into actionable 
information.

\subsection{Improving Existing Code Structure}
Another important direction concerns the maintainability of the implementation. The current code validates the main 
concepts of the thesis, but future versions should make the pipeline more modular. Data extraction, preprocessing, network 
generation, platform matching, routing, trip generation, simulation execution, and output analysis could be separated into 
independent modules with clearer interfaces.

This would make the framework easier to test, extend, and reuse. It would also allow each component to be evaluated 
independently. For example, platform matching could be tested without running a full simulation, and routing algorithms 
could be compared on the same infrastructure graph. Adding automated validation tests would also help detect regressions 
when modifying the pipeline.

\subsection{Counterfactual Simulation and What-If Scenarios}
A further direction concerns counterfactual simulation and what-if scenarios. The current framework mainly reproduces 
existing infrastructure and train movements. A future version could allow users to modify parts of the network and observe 
the impact of these changes on railway traffic.

For example, it could be possible to simulate the closure of a station, the removal of a track section, a modification of 
train frequency, or a change in departure times. At the macro level, this could be implemented by modifying 
station-to-station connections. At the micro level, specific platforms, switches, or track segments could be made 
unavailable. This would make the digital twin useful as an exploratory tool for testing operational decisions before 
applying them to the real network.

\subsection{Tracks Works Planning}
Track works planning is another promising extension. Maintenance operations often require temporary closures, reduced 
capacity, speed restrictions, or rerouting. A simulation-based digital twin could help evaluate the impact of these 
interventions before they take place.

In the current framework, this could be implemented by adding availability constraints to infrastructure elements. At the 
macro level, railway connections could be removed or penalized. At the micro level, specific tracks, platforms, or switches 
could be marked as unavailable during a given time window. This would help compare maintenance strategies, identify 
bottlenecks, and reduce the operational impact of planned works.

\subsection{Real-Time Forecasting and Incident Propagation}
Another future direction is the integration of real-time forecasting and incident propagation. The current implementation 
is mainly based on historical and static data. A future version could integrate live or near-real-time information, such as 
train delays or incidents, and use it to estimate the future state of the network.

This extension could combine the simulation pipeline with statistical or machine learning models trained on punctuality 
data. The macro-level model could be used to study delay propagation between stations, while the micro-level model could 
help analyse the influence of local infrastructure constraints. Such an extension would move the framework closer to an 
operational digital twin for monitoring and proactive traffic management.

\subsection{Framework Generalization}
A final direction concerns the generalization of the framework. The current implementation is adapted to the Belgian 
railway datasets used in this thesis. A more general version should separate the core modeling logic from dataset-specific 
preprocessing steps.

This could be achieved by introducing modular data adapters that convert different railway datasets into a common internal 
representation. Depending on the available data, the framework could then generate a macro-level model, a micro-level 
model, or a hybrid representation. This would make the approach easier to transfer to other railway networks and would 
improve its long-term usability.

\section{Final Remarks}
This thesis shows that building a railway digital twin is not only a simulation problem, but also a data integration and 
infrastructure modelling problem. The results demonstrate that a complete pipeline can be constructed from real railway 
datasets and executed in SUMO, especially at the macro level. They also show that moving toward a detailed micro-level 
representation introduces new challenges, particularly in route reconstruction and topological consistency.

The most important lesson of this work is that infrastructure coverage alone is not sufficient. A generated railway network 
must also be operationally usable: platforms must be reachable, track directions must be coherent, switches must preserve 
connectivity, and real train paths must be routable. The failure of the real-data micro-level simulation is therefore not 
only a limitation, but also a valuable finding. It identifies the next technical obstacle to solve in order to move from a 
structural digital representation toward a more complete operational railway digital twin.

Overall, this work provides a first foundation for simulation-oriented railway digital twin development. It validates the 
feasibility of the macro-level approach, demonstrates the potential of the micro-level approach, and identifies the 
improvements required to make future versions more realistic, robust, and useful for railway analysis.